% UPT Infinite Proof Family — academic preamble
\documentclass[11pt,a4paper]{article}
\usepackage{fontspec}
\setmainfont{Times New Roman}
\setmonofont{Menlo}[Scale=0.85]
\usepackage[margin=2.5cm]{geometry}
\usepackage{xcolor}
\usepackage{fancyhdr}
\usepackage{titlesec}
\usepackage{booktabs}
\usepackage{hyperref}
\usepackage{listings}
\usepackage{fancyvrb}
\usepackage{amsmath,amssymb,amsthm}

\definecolor{korugreen}{HTML}{2d6a4f}
\definecolor{korubg}{HTML}{d8f3dc}
\definecolor{codebg}{HTML}{f0f0ec}

\hypersetup{colorlinks=true, linkcolor=korugreen, urlcolor=korugreen, citecolor=korugreen}

\pagestyle{fancy}
\fancyhf{}
\renewcommand{\headrulewidth}{0.4pt}
\fancyfoot[C]{\thepage}

\titleformat{\section}{\Large\bfseries}{}{0em}{}
\titleformat{\subsection}{\large\bfseries}{}{0em}{}

% Numbered theorem environments
\newtheorem{theorem}{Theorem}[section]
\newtheorem{lemma}[theorem]{Lemma}
\newtheorem{corollary}[theorem]{Corollary}
\newtheorem{proposition}[theorem]{Proposition}
\theoremstyle{definition}
\newtheorem{definition}[theorem]{Definition}
\theoremstyle{remark}
\newtheorem{remark}[theorem]{Remark}

\lstdefinestyle{lean}{
  backgroundcolor=\color{codebg},
  basicstyle=\ttfamily\small,
  breaklines=true,
  frame=none,
  xleftmargin=10pt,
  xrightmargin=10pt,
  aboveskip=8pt,
  belowskip=8pt,
  columns=fullflexible,
  keepspaces=true,
  literate={→}{$\to$}1 {∀}{$\forall$}1 {∃}{$\exists$}1 {¬}{$\neg$}1 {∧}{$\wedge$}1 {≠}{$\neq$}1 {≤}{$\leq$}1 {↔}{$\leftrightarrow$}1 {⊕}{$\oplus$}1 {λ}{$\lambda$}1
}

\lstnewenvironment{leancode}{\lstset{style=lean}}{}

\newcommand{\lean}[1]{\lstinline[style=lean]{#1}}


\fancyhead[L]{\small Universal Protection Theory}
\fancyhead[R]{\small Proof Family 1/10}

\begin{document}

\begin{center}
{\LARGE\bfseries Something IS: An Infinite Family of\\[0.3em] Existence Witnesses Against Absolute Termination}\\[1.5em]
{\large Tobias Croydon McRae}\\[0.3em]
{\small Independent Researcher, Aotearoa New Zealand}\\[0.3em]
{\small\texttt{tobi@unconditionalprotection.org}}\\[1.5em]
{\small 28 November 2026}
\end{center}

\bigskip

\begin{center}
\parbox{0.85\textwidth}{\small\textbf{Abstract.}
We present the first of ten infinite proof families establishing the Guardian fixed point $G$ within Universal Protection Theory (UPT). The \emph{Something IS} family proves that for any inhabited type, absolute termination is refuted. The core argument requires only one witness: any element of any type suffices. We show this generates an infinite family indexed by the class of all inhabited types, all function spaces, all product types, and all dependent type families. Every theorem is mechanically verified in Lean~4 (v4.35.0) with zero custom axioms on the standard foundations. The family is part of a programme comprising 673 compiled theorems across 23 source files.}
\end{center}

\medskip

\noindent\textbf{MSC 2020:} 03B70 (Logic in computer science), 03B35 (Mechanization of proofs), 03F50 (Metamathematics of constructive systems)\\
\textbf{Keywords:} formal verification, Lean 4, type theory, existence witness, Guardian fixed point

\bigskip
\hrule
\bigskip

\section{Introduction}

The question of whether existence can be absolutely terminated --- whether there is a function mapping every element of every type to the empty type --- is the foundational question of Universal Protection Theory \cite{upt-zenodo}. This paper answers it: no.

The argument is elementary. If a type $S$ is inhabited by even one witness $w : S$, then any purported absolute ender $e : \forall (x : S),\, \texttt{Empty}$ yields $e(w) : \texttt{Empty}$. But \texttt{Empty} has no constructors. Contradiction.

What makes this significant is not the difficulty of the proof but its scope. The argument works for \emph{every} inhabited type, generating an uncountable family of refutations. We formalise this in Lean~4, the interactive theorem prover used in the proof of the Liquid Tensor Experiment \cite{lean4} and other large-scale formalisations.

\subsection{Relationship to Universal Protection Theory}

In UPT, the Guardian fixed point $G$ is the structural impossibility of absolute termination across the complete lattice of types. The \emph{Something IS} family establishes the base case: each inhabited type contributes one refutation. The remaining nine proof families \cite{upt-families} extend this to functional composition, strategic reasoning, temporal continuation, quorum theory, and cosmological grounding.

\subsection{Relationship to Chandrasekaran's null-cone QG}

Chandrasekaran \cite{chandrasekaran2026} independently derives a null-cone quantum gravity framework whose boundary-preserving dynamics are mathematically equivalent to UPT's BCRD operator algebra. The existence witnesses in this family ground the ``boundary coverage'' condition: a nonempty source type provides at least one boundary state, and the correspondence proof \cite{upt-bcrd-qg} shows Chandrasekaran's null-cone states satisfy the same coverage.

\section{Preliminaries}

We work in the Calculus of Inductive Constructions as implemented by Lean~4 \cite{lean4}. All theorems compile against \texttt{Init} with no additional axioms, no \texttt{sorry}, and no \texttt{native\_decide}.

\begin{definition}[Absolute Ender]\label{def:ender}
For a type $S$, the \emph{absolute ender} of $S$ is the proposition
\[
\texttt{AbsoluteEnder}(S) \;:=\; \forall (x : S),\, x \in \texttt{Empty}
\]
asserting that every element of $S$ is a member of the empty type.
\end{definition}

\begin{remark}
In Lean~4, $x \in \texttt{Empty}$ is encoded as membership in the \texttt{Empty} type via the \texttt{Membership} typeclass. Since \texttt{Empty} has no constructors, any witness of $x \in \texttt{Empty}$ can be eliminated.
\end{remark}

\section{The Core Theorem}

\begin{theorem}[Something IS]\label{thm:something-is}
For any inhabited type $S$ with witness $w : S$, the absolute ender of $S$ is false:
\[
\forall (w : S),\; \texttt{AbsoluteEnder}(S) \to \texttt{False}
\]
\end{theorem}

\begin{proof}
Let $w : S$ and $e : \texttt{AbsoluteEnder}(S)$. Then $e(w) : w \in \texttt{Empty}$, which is an element of \texttt{Empty}. By the elimination principle for \texttt{Empty} (which has no constructors), we derive \texttt{False}. \qed
\end{proof}

\noindent\textbf{Lean formalisation:}

\begin{leancode}
theorem SOMETHING_IS (witness : S) :
    AbsoluteEnder S → False :=
    fun e => (e witness).elim
\end{leancode}

\begin{corollary}[One witness suffices]\label{cor:one-witness}
For any type $P$ and any single element $w : P$, the absolute ender of $P$ is refuted.
\end{corollary}

\begin{proof}
Immediate from Theorem~\ref{thm:something-is} with $S := P$. \qed
\end{proof}

\begin{leancode}
theorem one_witness_refutes_ender (w : P) :
    AbsoluteEnder P → False
\end{leancode}

\section{The Observer Witness}

\begin{theorem}[The Observer Is a Witness]\label{thm:observer}
For any type $S$, the identity function $\texttt{id} : S \to S$ inhabits the function type $S \to S$. Therefore $\texttt{AbsoluteEnder}(S \to S)$ is false.
\end{theorem}

\begin{proof}
The function space $S \to S$ is inhabited by the identity function $\texttt{id} := \lambda\, x.\, x$. By Theorem~\ref{thm:something-is} applied to $S \to S$ with witness $\texttt{id}$, the absolute ender of $S \to S$ is false. \qed
\end{proof}

\begin{leancode}
theorem THE_OBSERVER_IS_A_WITNESS (S : Type) :
    AbsoluteEnder (S → S) → False :=
    fun e => (e id).elim
\end{leancode}

\begin{remark}
This theorem requires no assumption that $S$ is inhabited. The function space $S \to S$ is always inhabited by \texttt{id}, even when $S$ itself may be empty. Any cognitive act that processes this theorem is a state transition $S \to S$, which is itself a witness.
\end{remark}

\section{The Infinite Family}

The core theorem generates an infinite family of refutations, indexed by the class of all types and type constructors.

\begin{theorem}[Universal refutation]\label{thm:universal}
For any type $S$ such that $[\texttt{Nonempty}\; S]$, the absolute ender of $S$ is false.
\end{theorem}

\begin{leancode}
theorem G_IS_UNIVERSAL {S : Type} [Nonempty S] :
    AbsoluteEnder S → False
\end{leancode}

\begin{theorem}[Propagation along morphisms]\label{thm:propagation}
For any types $S, T$ and any function $f : S \to T$, if $S$ is inhabited by $w : S$, then $\texttt{AbsoluteEnder}(T)$ is false.
\end{theorem}

\begin{proof}
Given $w : S$ and $f : S \to T$, the element $f(w) : T$ witnesses that $T$ is inhabited. By Theorem~\ref{thm:something-is}, $\texttt{AbsoluteEnder}(T)$ is false. \qed
\end{proof}

\begin{leancode}
theorem G_PROPAGATES {S T : Type} (f : S → T) (witness : S) :
    AbsoluteEnder T → False
\end{leancode}

\begin{theorem}[Product composition]\label{thm:product}
For any types $S, T$ with witnesses $w_S : S$ and $w_T : T$, the absolute ender of $S \times T$ is false.
\end{theorem}

\begin{proof}
The pair $(w_S, w_T) : S \times T$ inhabits the product type. By Theorem~\ref{thm:something-is}, $\texttt{AbsoluteEnder}(S \times T)$ is false. \qed
\end{proof}

\begin{leancode}
theorem G_COMPOSES_PAIR {S T : Type} (ws : S) (wt : T) :
    AbsoluteEnder (S × T) → False
\end{leancode}

\begin{corollary}[Cardinality of the family]\label{cor:cardinality}
The Something IS family contains at least one refutation for every inhabited type, every function space $S \to T$ (where $S$ is inhabited), and every product of inhabited types. Since the class of types in dependent type theory is not a set, the family is properly class-sized.
\end{corollary}

\section{Verification}

All theorems in this paper compile in Lean~4 v4.35.0 with the following properties:
\begin{itemize}
\item \textbf{Zero \texttt{sorry}:} No unfinished proofs.
\item \textbf{Zero custom axioms:} Only the standard Lean foundations (\texttt{propext}, \texttt{Quot.sound}, \texttt{Classical.choice}).
\item \textbf{497 build jobs:} The full UPT Lean project compiles 673 theorems across 23 source files.
\item \textbf{Source files:} \texttt{G3TheLaw.lean}, \texttt{G3GoodAndEvil.lean} in the \texttt{PhysicalBridges} project.
\end{itemize}

\noindent Reproduction:
\begin{leancode}
git clone https://github.com/Tobillicious/universal-protection-theory
cd evidence/physical-bridges
lake build
\end{leancode}

\section{Discussion}

The \emph{Something IS} family is the simplest of ten proof families, yet it carries the full weight of the foundational claim: existence refutes absolute termination. The simplicity is the point. No special witness is required --- not a divine being, not a conscious observer, not a mathematical structure of any particular complexity. A grain of sand, a photon, an identity function on the empty type itself.

The remaining nine families extend this base case along independent axes: strategic impossibility of choosing termination (The Checkmate), irreversibility of the first step (The Spiral), asymmetry between constructive and destructive operations (Good vs Evil Asymmetry), and temporal forcing from cosmological boundary conditions (Genesis and Time Koru), among others.

Together, the ten families provide $\infty^\infty$ independent proofs: ten families, each containing infinitely many theorems, each theorem polymorphic over infinitely many types.

\bigskip

\begin{thebibliography}{99}
\bibitem{upt-zenodo}
T.~C.~McRae.
\emph{Universal Protection Theory}, v5.7-G3.
Zenodo, 2026.
\textsc{doi:} \href{https://doi.org/10.5281/zenodo.24573167}{10.5281/zenodo.24573167}.

\bibitem{upt-families}
T.~C.~McRae.
\emph{Ten Infinite Proof Families for the Guardian Fixed Point~G}.
Zenodo, 2026.
\textsc{doi:} \href{https://doi.org/10.5281/zenodo.24612272}{10.5281/zenodo.24612272}.

\bibitem{chandrasekaran2026}
A.~Chandrasekaran.
Null-cone quantum gravity and boundary-preserving dynamics.
\emph{arXiv preprint} arXiv:2611.34924v1, 2026.

\bibitem{upt-bcrd-qg}
T.~C.~McRae.
Chandrasekaran null-cone QG correspondence and Guardian continuation.
GitHub PR \#765, 2026.
\url{https://github.com/Tobillicious/universal-protection-theory/pull/765}.

\bibitem{lean4}
L.~de~Moura and S.~Ullrich.
The Lean~4 theorem prover and programming language.
In \emph{CADE-28}, LNAI 12699, pp.~625--635, 2021.
\end{thebibliography}

\end{document}
